\documentclass[11pt]{article}
\usepackage[a4paper,margin=25mm]{geometry}
\usepackage{amsmath,amssymb,amsthm,booktabs,hyperref,xurl}
\newtheorem{theorem}{Theorem}
\newtheorem{lemma}[theorem]{Lemma}
\newtheorem{proposition}[theorem]{Proposition}
\newcommand{\HH}{\mathbb H}
\title{Forced mixed corners and low-facet rigidity\linebreak for right-angled hyperbolic polytopes}
\author{Research manuscript with reproducible exact calculations}
\date{18 October 2026}
\begin{document}
\maketitle
\begin{abstract}
Any finite-volume right-angled hyperbolic polytope in dimension eleven
or twelve contains a specific ten-facet mixed four-face. We derive
explicit finite/ideal vertex ratios and positive occurrence bounds
using a weighted two-face inequality. We also classify equality in the
four-dimensional facet-plus-cusp bound fifteen and in the
five-dimensional facet bound sixteen. A complete five-dimensional
corner forces either that sixteen-facet polytope or at least twenty
facets; in every other five-polytope such corners occupy at most one
eleventh of the finite vertices. These are universal necessary
conditions, with no arithmetic or facet-spectrum restriction. They do
not improve the published dimension ceiling twelve. Exact rational
calculations, finite model checks and a numerical feasibility diagnostic
are supplied; the geometric proofs have not yet been formalized in Lean.
\end{abstract}

\section{Statements and conventions}
Let $P\subset\HH^n$ be a finite-volume right-angled polytope. Write
$f,c,a_j$ for its finite vertices, ideal vertices and $j$-faces.
Faces of positive dimension inherit right angles and finite volume.
At finite vertices the links are simplices; at ideal vertices they are
rectangular boxes. Outward facet normals have square one, and the
Lorentz metric has one negative direction.

Let $P_3$ be the six-facet triangular bipyramid with two finite and
three ideal vertices. Let $P_4$ be the ten-facet polytope with five
finite and five ideal vertices described in Section~\ref{sec:corners}.
Let $P_5$ be the sixteen-facet polytope with sixteen finite and ten
ideal vertices described in Section~\ref{sec:p5}.
These are the known Potyagailo--Vinberg polytopes~\cite{PV}.
Let $T$ count triangular two-faces with one finite vertex, and let
$B,C,D$ count faces isometric to $P_3,P_4,P_5$, respectively.

\begin{theorem}\label{thm:density}
For $n=12$,
\[
 \frac{1474}{21}c<f<\frac{1298}{9}c,\qquad
 T>\frac{331}{437}a_2,\qquad T>\frac{1986}{59}f,
 \qquad C>\frac{627}{1475}f.
\]
For $n=11$,
\[
 \frac{268}{7}c<f<\frac{236}{3}c,\qquad
 T>\frac{331}{437}a_2,\qquad T>\frac{1655}{59}f,
 \qquad C>\frac{1763}{5310}f.
\]
In either dimension $P$ has a four-face isometric to $P_4$.
\end{theorem}

\begin{theorem}\label{thm:rigid}
A right-angled four-polytope with facet-plus-cusp sum fifteen is
isometric to $P_4$. A right-angled five-polytope with sixteen facets
is isometric to $P_5$. Consequently
\[
 C\ge 11a_4-(n-3)a_3,\qquad
 D\ge 17a_5-(n-4)a_4.
\]
\end{theorem}

Our published inputs are the standard truncation and
Dehn--Sommerville framework described in~\cite{Alex}, the lower bound
eighty-two on eight-dimensional facet counts recorded there, and the
four-dimensional bounds $F\ge10$, $F+c\ge15$ in
\cite[Section~4, equations~(7)--(8)]{PV}. Compact right-angled polytopes
have dimension at most four~\cite{PV}, so the noncompact eight-dimensional
bound applies to every eight-face here. Equality classification is
proved below; it is not assumed from published attainment of the bounds.

\section{The weighted two-face calculation}
Put $\lambda=2123/1280$. For $n=12$, truncation at the cusps gives
\begin{align*}
 f&=a_6-5a_7+14a_8-28a_9+42a_{10}-42a_{11}-462c,\\
 a_2&=15a_6-74a_7+206a_8-411a_9+616a_{10}-616a_{11}-2684c.
\end{align*}
These rows follow directly by setting $b_0=f+2048c$, $b_{11}=a_{11}+c$
and $b_k=a_k+2^{11-k}\binom{11}{k}c$ for $1\le k\le10$, and equating
the coefficients of $H(t)$ and $t^{12}H(1/t)$ for
\[
 H(t)=(t-1)^{12}+\sum_{k=0}^{11}b_k(t-1)^k.
\]
Eliminating gives
\begin{equation}\label{eq:weighted}
 66f+11264\lambda c-\frac{22}{5}a_2
 =\frac{-22a_7+88a_8-198a_9+308a_{10}-308a_{11}}5<0.
\end{equation}
Indeed $5a_7\ge82a_8$ and $3a_9\ge11a_{10}$, so in the numerator
the first pair of terms is at most $-1364a_8/5$ and the second at most
$-418a_{10}$.

A polygon with $u$ finite and $v$ ideal vertices satisfies
$u+v\ge3$ and $u+2v\ge5$, the latter by positive hyperbolic area.
Summing these two inequalities and using~\eqref{eq:weighted} yields
the two vertex-ratio bounds in dimension twelve. A counted triangle
has weighted size $1+2\lambda$; every other polygon has size at least
$3+\lambda$. Hence
\[
 (3+\lambda)a_2-(2-\lambda)T
 \le66f+11264\lambda c<\frac{22}{5}a_2,
\]
which gives $T>331a_2/437$.

For every noncounted polygon,
\[
 u+\lambda v-\frac{22}{5}\ge\frac{331}{3840}u.
\]
It suffices to check the minimal pairs $(0,3),(2,2),(3,1),(5,0)$;
every other permissible noncounted pair dominates one of them.
A counted triangle has deficit $106/1280$.
Summing gives $331(66f-T)<318T$, and therefore $T>1986f/59$.

For dimension eleven the exact coefficient rows are
\begin{align*}
 f&=2a_6-12a_7+42a_8-112a_9+252a_{10}-252c-504,\\
 a_2&=25a_6-148a_7+515a_8-1370a_9+3080a_{10}-1220c-6160,\\
 55f+5120\lambda c-\frac{22}{5}a_2
 &=-\frac{44}{5}a_7+44a_8-132a_9+308a_{10}-616<0.
\end{align*}
Here $4a_7\ge82a_8$ and $2a_9\ge11a_{10}$ prove the strict sign.
The same polygon argument gives the stated eleven-dimensional ratios
and triangle densities.

\section{Actual three- and four-faces}\label{sec:corners}
At a finite vertex let $q^2=-1$ and let $n_1,\ldots,n_n$ be its
orthonormal outward normals. Form a graph $G_q$ on the edge directions,
joining $i,j$ when their two-face is counted by $T$. Its opposite
unit wall is uniquely
\begin{equation}\label{eq:m}
 m_{ij}=q-n_i-n_j.
\end{equation}
Its ideal endpoints are $q-n_i,q-n_j$.

\begin{lemma}\label{lem:p3}
A triangle in $G_q$ supplies an actual three-face isometric to $P_3$.
\end{lemma}
\begin{proof}
For the triangle $i,j,k$, the opposite point is
$q'=2q-n_i-n_j-n_k$, of square $-1$.
For any other wall write $r=a q-\sum c_l n_l$, where $a>0$ and
$c_l\ge0$. Acuteness with the three normals~\eqref{eq:m} gives
$c_i+c_j,c_i+c_k,c_j+c_k\le a$. Thus
$r\cdot q'=-2a+c_i+c_j+c_k\le-a/2<0$.
The selected three $m$ normals and the $n_l$ outside the triple vanish
there and are orthonormal. The six walls in the resulting three-space
bound the convex hull of $q,q'$ and the three ideal endpoints.
All five points are feasible, so convexity identifies the whole face.
\end{proof}

Every right-angled three-polytope has at least six facets, with equality
only for $P_3$. To prove this, write $p,s,F$ for finite vertices, cusps
and facets. Euler and vertex degrees give $p+2s=2F-4$; summed polygon
areas give $p+6s\ge20$. Together these exclude $F\le5$.
If $F=6$, they give $s\ge3$. The polygon-corner sum is $24-2s$, so
at least $2s$ of the six facets are triangular, forcing $s\le3$.
Thus $s=3,p=2$ and all six polygons are counted triangles. Lemma~\ref{lem:p3}
identifies the polytope.

Double counting now gives
\[
 \sum_q e(G_q)=T,\qquad \sum_q t(G_q)=2B.
\]
In dimension twelve, three-face Euler and the six-facet gap give
\[
 10a_2=2a_3+110f+28160c,\qquad B\ge7a_3-10a_2.
\]
Since~\eqref{eq:weighted} gives $a_2>15f+4246c$, we obtain
\[
 B>7590c-10f>\frac{2515}{59}f.
\]
A four-clique-free graph on twelve vertices has at most sixty-four
triangles. One elementary proof maximizes the weighted triangle sum
on the vertex-weight simplex with minimum support: two nonadjacent
supported vertices can be merged, so the support is a clique of size
at most three. Its product is at most $1/27$; uniform weights give
$12^3/27=64$. Thus some $G_q$ has a four-clique.

For such a four-set $S$, its four $n_i$ and six $m_{ij}$ define the
projective region $y_i\ge0$, $y_i+y_j\le t$, with
$x=tq-\sum_{i\in S}y_i n_i$. Its vertices are $q$, the four points
$2q-\sum_{i\in A}n_i$ for three-subsets $A\subset S$, four ideal
points $q-n_i$, and the ideal point $2q-\sum_{i\in S}n_i$.
The last is $m_{ij}+m_{kl}$ for a partition of $S$, and has nonpositive
product with every other wall by acuteness. The remaining points are
feasible by Lemma~\ref{lem:p3} and the original triangles.
Their finite-volume convex hull is therefore the entire four-face:
it is $P_4$.

Every four-polytope has at least ten facets. Indeed the preceding
three-dimensional inequalities give $F+s\ge9$ for each three-facet;
its distinct neighbors and one distinct cusp-opposite per cusp force
at least $1+F+s$ ambient facets. At equality every facet has $F=6,s=3$,
so the preceding four-clique argument identifies the whole polytope
as $P_4$. Every other four-face has at least eleven facets, proving
$C\ge11a_4-(n-3)a_3$.

Deleting an edge of a four-clique until none remains takes at most
$k_4(G)$ deletions. Each loses at most ten triangles on twelve vertices,
so $t(G)\le64+10k_4(G)$. As $\sum_q k_4(G_q)=5C$, the lower bound on
$B$ gives $C>627f/1475$.
In dimension eleven the corresponding identities are
\[
 a_3=\frac92a_2-\frac{165}{4}f-5760c,\qquad
 B>3105c-\frac{15}{2}f>\frac{7545}{236}f.
\]
The triangle bound is $11^3/27$, with integer part forty-nine;
edge deletion loses at most nine triangles per edge.
This yields $C>1763f/5310$ and completes Theorem~\ref{thm:density}.

\section{Minimum-sum four-dimensional rigidity}
Distinct outward unit normals have product zero or at most $-1$.
For if $-1<u\cdot v<0$, their positive definite Gram matrix has
nonnegative inverse. Projecting a strict timelike center onto their
orthogonal complement preserves all other strict wall inequalities
by acuteness and gives an actual nonright ridge, a contradiction.

\begin{lemma}\label{lem:two}
Two counted triangular facets sharing their ideal--ideal edge in a
right-angled three-polytope force the whole polytope to be $P_3$.
\end{lemma}
\begin{proof}
Normalize the common null endpoints by $u\cdot v=-1$ and let the
orthogonal facet normals be $n_1,n_2$. Their finite apices are
$u+v-n_2$ and $s u+s^{-1}v-n_1$, $s>0$. The other four walls have
normals $u-n_2,v-n_2,su-n_1,s^{-1}v-n_1$. Two cross products are
$-s,-1/s$, so the normal-product gap forces $s=1$.
The third cusp $u+v-n_1-n_2=(u-n_1)+(v-n_2)$ is feasible by acuteness.
All five vertices of the six-wall bipyramid are therefore feasible,
proving the claim by convexity.
\end{proof}

Suppose a four-polytope has $F+c=15$. The published $F\ge10$ gives
$c\le5$. A three-facet with $r$ polygon facets and $s$ cusps satisfies
$r+2s\ge12$. Summing gives degree sum at least $12F-12c$ in the facet
graph, since every cusp lies in six facets. This graph has no five-clique:
mutually adjacent normals would be five orthogonal spacelike vectors
in positive index four. Its degree sum is at most $3F^2/4$ by the
elementary weighted clique bound. The lower bound exceeds this for
$c=0,1,2$. For $c=3,F=12$, equality forces all degrees nine: uniform
weights maximize the edge quadratic in the interior of the simplex,
so its first derivatives are equal. But equality also gives
$r+2s=12$ in every facet, making every degree even, a contradiction.

For $c=4,F=11$, the neighbor/cusp-opposite injection gives $r+s\le10$.
Together with $r+2s\ge12$ this gives $s\ge2$. Values $s\ge4$ would
give $r\le6$, contrary to the six-facet equality classification.
Since $\sum s=24$, nine facets have two cusps and two have three.
A two-cusp facet has $r=8$ and zero total polygon excess
$(\hbox{sides})+(\hbox{ideal vertices})-5$. If it has a triangle, the
polygon across its ideal--ideal edge is also a counted triangle;
Lemma~\ref{lem:two} contradicts its eight facets.
Its thirty-two polygon corners thus make all eight polygons
quadrilaterals. A three-cusp facet has six or seven polygons, at least
four triangular. Each would need a different adjacent facet with a
triangular polygon, whereas only the other three-cusp facet can do so.
This excludes $c=4$. Finally $c=5$ gives $F=10$ and hence $P_4$.

\section{The sixteen-facet five-polytope}\label{sec:p5}
The injection and $F+c\ge15$ force at least sixteen facets in dimension
five. At equality every four-facet has sum fifteen and is $P_4$.
Pick a finite vertex $q$ with orthonormal normals $n_1,\ldots,n_5$.
All ten $m_{ij}$ of~\eqref{eq:m} are actual walls, giving fifteen walls.
For each four-set $S$ the forced $P_4$ has cusp
$v_S=2q-\sum_{i\in S}n_i$. Exactly seven known walls meet it: the
omitted $n_k$ and six $m_{ij}$. Its rectangular four-box link needs
one more wall, the remaining sixteenth wall $r$, opposite $n_k$.
Thus $r+n_k=L_kv_S$, $L_k>0$. Comparing the $q$ coefficient and then
an $n_k$ coefficient makes every $L_k=1$. Consequently
\[
 r=2q-\sum_{i=1}^5n_i.
\]
The sixteen walls define $y_i\ge0$, $y_i+y_j\le1$, $\sum y_i\le2$.
Its exact vertex enumeration gives the origin; five ideal axes; ten
finite half-triples; five ideal half-four-sets; and five finite points
with one coordinate $2/3$ and the others $1/3$.
There are sixteen finite and ten ideal vertices. Every facet is $P_4$.
Thus the polytope is $P_5$, proving Theorem~\ref{thm:rigid}.

\section{Complete five-corners and a sharp twenty-facet gap}
Call a finite corner complete when its triangle graph is complete.
\begin{theorem}\label{thm:twenty}
A right-angled five-polytope with a complete corner is either $P_5$
or has at least twenty facets. In the latter case, if $L,V$ count its
complete corners and finite vertices, then $11L\le V$.
Both bounds are sharp for the reflected double of $P_5$ across the
wall $\rho=2q-\sum_{i=1}^5n_i$: it has twenty facets, twenty-two finite
vertices, fifteen cusps and two complete corners.
\end{theorem}
\begin{proof}
The five $n_i$ and ten $m_{ij}$ at a complete corner supply fifteen
actual walls. The canonical sixteen-wall $P_5$ seed is contained in
the polytope even if $\rho$ is absent. All but its five extra finite
vertices are feasible by the three- and four-face constructions.
For the remaining points $z_i=3q-2n_i-\sum_{j\ne i}n_j$, write an
additional wall as $r=a q-\sum c_j n_j$. We need
$2c_i+\sum_{j\ne i}c_j\le3a$, knowing $c_j\ge0$ and
$c_i+c_j\le a$. If $c_i\ge a/2$, summing its four pair bounds gives
at most $4a-2c_i\le3a$. Otherwise split the other four indices into
two pairs to get at most $2c_i+2a\le3a$. Convexity proves containment.

At each cusp $v_i=2q-\sum_{j\ne i}n_j$, the six $m_{jk}$ form three
opposite pairs; with $n_i$ these are seven of its eight walls.
The remaining wall is $r_i=\lambda_i v_i-n_i$, $\lambda_i>0$.
Since $r_i\cdot m_{ij}=1-\lambda_i$, the normal-product gap forces
$\lambda_i=1$ or $\lambda_i\ge2$.
For different $i,j$,
\[
 r_i\cdot r_j=1-(\lambda_i-1)(\lambda_j-1).
\]
If some $\lambda_i=1$, product one forces all these walls to coincide;
coefficient comparison makes every $\lambda_j=1$ and the common wall
is $\rho$. The sixteen actual walls and seed containment give $P=P_5$.
Otherwise the five $r_i$ are distinct and have $q$ coefficient at least
four, so none is among the first fifteen. There are at least twenty walls.

Suppose now $P\ne P_5$. Each complete corner has ten distinct opposite
finite vertices across its $P_3$ faces. Such a recipient is not complete:
its edge on the two untouched normals would otherwise force $\rho$.
In fact a recipient admits at most one complete predecessor.
To see this, label a receiving triple $S=\{1,2,3\}$ at a recipient $p$.
The complete opposite is $p_S=2p-n_1-n_2-n_3$, with incident normals
$m_{23},m_{13},m_{12},n_4,n_5$. Its complete graph forces all cross walls
\[
 p_S-m_{ij}-n_l=p-n_k-n_l,\qquad \{i,j,k\}=S,\quad l\notin S.
\]
These supply actual counted triangles at $p$: their null endpoints are
sums of known outward normals and hence feasible, and the selected
three walls bound their entire triangle. Therefore the receiving graph
is exactly $K_5$ minus the untouched edge. Its unique missing edge fixes
$S$ and the predecessor. Counting gives $10L\le V-L$.
The exact reflected-double calculation proves sharpness.
\end{proof}

If $K=\sum_q k_5(G_q)$ in ambient dimension $n$, summing over five-faces
gives the further universal inequality
\begin{equation}\label{eq:k5}
 11K\le\binom n5 f+160D.
\end{equation}
Indeed a canonical five-face has sixteen complete finite corners,
whereas every other five-face satisfies Theorem~\ref{thm:twenty}.

\section{Propagation cannot close in dimension at least nine}
\begin{proposition}
For $n\ge9$, if a finite-volume right-angled $n$-polytope has a complete
finite corner, some five-face with a complete corner is not $P_5$.
\end{proposition}
\begin{proof}
Across any $P_3$ opposite the vector
$h=3q-\sum_{i=1}^n n_i$ is unchanged: the three new normals are the
three $m_{ij}$, and the others stay fixed. Thus $h^2=n-9$ and
$h\cdot q=-3$ throughout each opposition component.
If every complete five-corner lies in $P_5$, a complete $n$-corner
propagates to every $P_3$ opposite. The three new directions and their
cross edges are forced by the old complete graph; every pair of
untouched directions is forced by its canonical five-face.
The resulting finite component is therefore regular of degree
$d=\binom n3$, entirely consisting of complete corners.
At each of its vertices,
\[
 \frac1d\sum_{|S|=3}q_S
 =\left(2-\frac9n\right)q+\frac3n h.
\]
Summing over its $N$ vertices yields
$(9-n)\sum q=3Nh$. For $n=9$ this contradicts $h\ne0$.
For $n>9$ it equates a timelike vector with a spacelike one.
\end{proof}
This is a propagation obstruction, not yet a quantitative density
contradiction. It does not by itself lower a dimension ceiling.

\section{A universal opposition moment bound in dimension twelve}
\begin{proposition}\label{prop:moment}
Every finite-volume right-angled twelve-polytope satisfies
\[
 2B\le\frac{1267552}{6875}f<184.372f.
\]
\end{proposition}
\begin{proof}
On each opposition component, $h^2=3$. Choose a future null vector
$z\perp h$, put $E(q)=-z\cdot q>0$ and $p_i=-z\cdot n_i/E(q)$.
Then $\sum p_i=3$, $\sum p_i^2=1$. For each of the 220 formal triples,
whether or not it supplies an actual face,
\[
 r_S=\frac{E(q_S)}{E(q)}=2-\sum_{i\in S}p_i.
\]
Set $x_i=p_i-1/4$, $\mu=5/4$ and $d_S=-\sum_{i\in S}x_i$.
Thus $r_S=\mu+d_S$, $\sum x_i=0$, $\sum x_i^2=1/4$.
Cauchy--Schwarz after subtracting the mean of the three-set indicator
gives $|d_S|\le3/4$, hence $1/2\le r_S\le2$.

Uniform averages over the triples satisfy
\begin{align*}
 \mathop{\rm av} d&=0,&
 \mathop{\rm av} d^2&=\frac9{176},\\
 \mathop{\rm av} d^3&=-\frac{27}{220}\sum x_i^3,&
 \mathop{\rm av} d^4&=\frac3{440}-\frac3{220}\sum x_i^4\le\frac3{440}.
\end{align*}
Expand the powers using the subset inclusion probabilities
\[
 1/4,\qquad 1/22,\qquad 1/220.
\]
Also
\[
 |\sum x_i^3|\le(\sum x_i^2)^{3/2}=1/8.
\]
The exact reciprocal expansion
\[
 \frac1{\mu+d}=\frac1\mu-\frac d{\mu^2}+\frac{d^2}{\mu^3}
 -\frac{d^3}{\mu^4}+\frac{d^4}{\mu^4(\mu+d)}
\]
therefore gives
\[
 \mathop{\rm av}\frac1{r_S}
 \le\frac45+\frac{64}{125}\frac9{176}
 +\frac{256}{625}\frac{27}{1760}
 +\frac{512}{625}\frac3{440}
 =\frac{28808}{34375}.
\]
With $g(q)=1/E(q)$, the sum of $g(q_S)/g(q)$ over actual neighbors is
at most 220 times this value. Summing over vertices, each undirected
opposition edge contributes $g(v)/g(u)+g(u)/g(v)\ge2$. There are $B$
edges, which proves the proposition.
\end{proof}

For linear objectives involving only clique counts on twelve vertices,
all corner graph vectors can be replaced by convex mixtures of the
77 complete multipartite vectors. Indeed, partition any maximizing graph
into maximal independent classes with identical external neighborhoods.
For two nonadjacent classes $A,B$, no clique meets both, so the objective
is a constant plus $|A|d_A+|B|d_B$. Replace the lower-contribution class
by clones of the other. Other classes remain intact; the class count
decreases and the objective does not decrease, even for signed objective
coefficients. Iteration gives a complete multipartite graph. Equality
of all linear support functions proves the convex-hull statement.
This does not identify an individual geometric corner graph as multipartite.

\section{Reproduction, diagnostic and verification scope}
The standard-library exact checkers include
\begin{quote}\small
\path{scripts/check_right_angled_mixed_weighted_triangles_research_20261018.py}\\
\path{scripts/check_right_angled_mixed_low_facet_rigidity_research_20261018.py}.
\end{quote}
The complete-corner and opposition-moment additions are checked by
\begin{quote}\small
\path{scripts/check_right_angled_mixed_complete_five_corner_research_20261018.py}\\
\path{scripts/check_right_angled_mixed_opposition_moments_research_20261018.py}.
\end{quote}
They recompute the Dehn--Sommerville equations by rational elimination,
the numerical inequalities, the twenty and 210 active-wall choices
for $P_3,P_4$, and all 4368 five-wall choices for $P_5$.
The additional checks reconstruct the twenty-facet double and its
corner graphs, and verify the exact moment coefficients and constants.
Generation, fresh replays and independent written/source reviews pass.
The scripts verify coefficient and model claims; the universal geometric
arguments are reviewed written proofs, not encoded by those enumerations.

The separate exact diagnostic
\begin{quote}\small
\path{scripts/check_right_angled_mixed_incidence_feasibility_research_20261018.py}
\end{quote}
constructs a simultaneous numerical solution of the earlier dimension-twelve
aggregate constraints, including canonical incidences and integer-scaled
Macaulay conditions. Its original assignment violates the new moment
cap. The moment packet supplies a modified assignment satisfying the
new scalar cap and packing inequality as well; its complete-corner
saturation still conflicts with the propagation argument above.
No quantitative inequality yet excludes every nearby assignment.
None of these assignments supplies hyperplanes, a boundary complex,
consistent cusp widths or geometric gluing.

Run \path{scripts/check_mixed_corner_package_20261018.py --build}
to replay the five small checkers sequentially and compile this paper.
The runner uses the shared exclusive computation lock and resource
gates, retains a compact audit and the PDF, and removes TeX intermediates.

No new example or dimension ceiling is claimed. The new geometric
theorems here await Lean verification. The repository's existing Lean
results do not certify these arguments merely because their final
inequalities are rational.

\begin{thebibliography}{9}
\bibitem{PV} L. Potyagailo and E. Vinberg,
On right-angled reflection groups in hyperbolic spaces,
\emph{Comment. Math. Helv.} 80 (2005), 63--73.
\url{https://ems.press/content/serial-article-files/43032}.
\bibitem{Alex} S. Alexandrov,
On ideal vertices of right-angled hyperbolic polyhedra,
2023, Sections 2--6.
\url{https://arxiv.org/html/2303.09533v1}.
\end{thebibliography}
\end{document}
